\documentclass[12pt]{article}
\usepackage[utf8]{inputenc}
\usepackage[T1]{fontenc}
\usepackage[margin=1in]{geometry}
\usepackage{setspace}
\usepackage{xcolor}
\usepackage[colorlinks=true, urlcolor=blue!60!black, linkcolor=blue!60!black]{hyperref}
\setlength{\parindent}{2em}
\setlength{\parskip}{1.25em}
\renewcommand{\baselinestretch}{1.0}
\title{\textit{Loper Bright Enterprises v. Raimondo}, 603 U.S. 369 (2024)}
\author{Case Brief}
\date{}
\begin{document}
\maketitle

\section{Case Metadata}

\begin{itemize}
  \item \textbf{Case Name:} \textit{Loper Bright Enterprises v. Raimondo} (consolidated with \textit{Relentless, Inc. v. Department of Commerce}, No. 22--1219)
  \item \textbf{Citation:} 603 U.S. 369 (2024) (\href{https://www.supremecourt.gov/opinions/23pdf/22-451_7m58.pdf}{SupremeCourt.gov Slip Opinion}; \href{https://www.law.cornell.edu/supremecourt/text/22-451}{Cornell LII}; \href{https://supreme.justia.com/cases/federal/us/603/22-451/}{Justia})
  \item \textbf{Court:} Supreme Court of the United States (on writs of certiorari to the D.C. Circuit, 45 F.4th 359 (2022), and the First Circuit, 62 F.4th 621 (2023))
  \item \textbf{Dates:} Argued January 17, 2024; Decided June 28, 2024
  \item \textbf{Vote:} 6--3 in \textit{Relentless} (No. 22--1219); 6--2 in \textit{Loper Bright} (No. 22--451, with Justice Jackson recused)
  \item \textbf{Majority Author:} Chief Justice John G. Roberts, Jr. (joined by Justices Thomas, Alito, Gorsuch, Kavanaugh, and Barrett)
  \item \textbf{Concurring Opinions:} Justice Thomas; Justice Gorsuch
  \item \textbf{Dissenting Opinion:} Justice Kagan (joined by Justice Sotomayor, and by Justice Jackson as to \textit{Relentless})
  \item \textbf{Chapter Placement:} Chapter 2: Separation of Powers and the Executive (Judicial Review of Agency Action; Overruling \textit{Chevron} Deference; Administrative Procedure Act \S~706)
\end{itemize}

\section{Facts and Procedural History}

To curb overfishing off the U.S. coast, Congress enacted the Magnuson-Stevens Fishery Conservation and Management Act (MSA), 16 U.S.C. \S~1801 \textit{et seq.}, administered by the National Marine Fisheries Service (NMFS) under the Secretary of Commerce. Under the MSA, regional fishery management plans may require domestic vessels to carry onboard federal ``observers'' to collect conservation data (\S~1853(b)(8)). The statute expressly authorizes industry-funded observer programs for three specific groups---foreign fishing vessels, vessels in limited access privilege programs, and vessels in the North Pacific Council (capping fees at 2--3\% of harvest value for the latter two)---but is silent on whether Atlantic herring fishermen must pay for observers (\href{https://www.law.cornell.edu/supremecourt/text/22-451}{603 U.S. at 380--381}).

In 2020, NMFS promulgated a final rule (85 Fed. Reg. 7414) approving an industry-funded monitoring program for the Atlantic herring fishery targeting 50\% trip coverage. Whenever NMFS required an observer but declined to supply a government-funded monitor, vessel owners had to contract with and pay certified third-party observers---estimated by NMFS at up to \$710 per day, reducing annual vessel returns by up to 20\% (\textit{id.} at 381).

Two sets of family-owned Atlantic herring fishing businesses sued under the MSA (which incorporates the Administrative Procedure Act (APA), 5 U.S.C. \S~551 \textit{et seq.}), arguing the MSA did not authorize NMFS to shift observer salaries onto domestic fishermen. In \textit{Loper Bright}, the District Court granted summary judgment to the Government (544 F. Supp. 3d 82 (D.D.C. 2021)), and a divided D.C. Circuit affirmed (45 F.4th 359 (2022)), finding the MSA not ``wholly unambiguous'' and deferring to NMFS as ``reasonable'' under Step Two of \href{https://supreme.justia.com/cases/federal/us/467/837/}{\textit{Chevron U.S.A. Inc. v. Natural Resources Defense Council, Inc.}}, 467 U.S. 837 (1984) (Judge Walker dissenting). In \textit{Relentless}, the District of Rhode Island (561 F. Supp. 3d 226 (2021)) and the First Circuit (62 F.4th 621 (2023)) likewise upheld the rule under \textit{Chevron}, though the First Circuit declined to specify whether its holding rested on Step One or Step Two. The Supreme Court granted certiorari in both cases limited to whether \textit{Chevron} should be overruled or clarified.

\section{Issue(s) Presented}

Should \textit{Chevron U.S.A. Inc. v. Natural Resources Defense Council, Inc.}, 467 U.S. 837 (1984)---which required federal courts to defer to ``permissible'' agency interpretations of ambiguous statutes---be overruled as inconsistent with the Administrative Procedure Act (5 U.S.C. \S~706) and the judicial function under Article III?

\section{Holding}

\textbf{Yes. \textit{Chevron} is overruled.} Vacating and remanding both judgments, the Court held that:

\begin{itemize}
  \item \textbf{Independent Judicial Judgment:} Under APA \S~706 and Article III's \textit{Marbury v. Madison} tradition, courts must exercise independent judgment in deciding whether an agency has acted within its statutory authority and may not defer to an agency interpretation merely because a statute is ambiguous (\href{https://www.law.cornell.edu/supremecourt/text/22-451}{603 U.S. at 412--413}).
  \item \textbf{\textit{Skidmore} Respect and Lawful Delegations Preserved:} Courts may still accord persuasive ``respect'' to an agency's thorough, consistent, and expert judgment under \textit{Skidmore v. Swift \& Co.}, 323 U.S. 134 (1944), and must respect statutes that expressly or constitutionally delegate discretionary authority (e.g., to ``fill up the details'' or regulate as ``appropriate'' or ``reasonable''), while policing statutory boundaries and reasoned decisionmaking (\textit{id.} at 394--396, 412--413).
  \item \textbf{Statutory \textit{Stare Decisis} Preserved:} Overruling \textit{Chevron} does not automatically disturb prior holdings that specific agency actions were lawful (including \textit{Chevron}'s Clean Air Act holding); mere prior reliance on \textit{Chevron} is not a ``special justification'' to overrule them (\textit{id.} at 412).
\end{itemize}

\section{Reasoning}

\subsection{Article III Tradition and APA \S~706 (Part II)}

Chief Justice Roberts grounded the Court's analysis in the Framers' design of Article III (The Federalist Nos. 37 and 78) and \textit{Marbury v. Madison}, 5 U.S. (1 Cranch) 137, 177 (1803): although statutory language is often obscure, resolving legal ambiguities is ``the proper and peculiar province of the courts'' (\href{https://www.law.cornell.edu/supremecourt/text/22-451}{603 U.S. at 385}). From the Founding through the New Deal, courts accorded ``respect''---especially to contemporaneous, consistent executive constructions (\textit{Edwards' Lessee v. Darby}; \textit{Skidmore})---and deferred to agency factfinding, but never surrendered independent judgment on pure questions of law. Even factbound New Deal decisions like \textit{Gray v. Powell} and \textit{NLRB v. Hearst Publications} reaffirmed that statutory interpretation remained for courts (\textit{id.} at 386--391).

Congress codified that tradition in the 1946 APA. Section 706 commands that ``the reviewing court shall decide all relevant questions of law, interpret constitutional and statutory provisions, and determine the meaning or applicability of the terms of an agency action'' (5 U.S.C. \S~706). Tellingly, while \S~706(2) prescribes deferential review for agency policymaking (``arbitrary, capricious'') and formal factfinding (``substantial evidence''), it prescribes no deference for legal questions---treating statutory and constitutional interpretation identically (\textit{id.} at 391--394).

\subsection{\textit{Chevron}'s Conflict with the APA (Part III)}

Decided by a six-Justice quorum without citing the APA, \textit{Chevron} inverted \S~706 by commanding courts to set aside their best reading of a statute in favor of an agency's ``permissible'' construction---even when the agency reversed course or overrode prior judicial precedent (\textit{Brand X}) (\textit{id.} at 396--399). Roberts dismantled \textit{Chevron}'s justifying presumption that statutory ambiguity equals an implicit congressional delegation to agencies: ambiguities are frequently unintentional, statutes possess a ``single, best meaning'' discernible through traditional interpretive tools, and agencies possess no special competence in legal construction (\textit{id.} at 399--401). Agency technical expertise can inform courts via \textit{Skidmore} without binding them; uniformity cannot justify adopting a wrong reading of a statute; and resolving legal ambiguity is judging, not policymaking. By forcing courts to treat ambiguity as delegation, ``\textit{Chevron} does not prevent judges from making policy. It prevents them from judging'' (\textit{id.} at 404).

\subsection{\textit{Stare Decisis} (Part IV)}

All three relevant \textit{stare decisis} factors favored overruling \textit{Chevron}: (1) \textbf{Quality of reasoning}---\textit{Chevron} ignored the APA and remained a ``rule in search of a justification''; (2) \textbf{Workability}---``ambiguity'' proved hopelessly indeterminate (``in the eye of the beholder''), spawning a ``dizzying breakdance'' of carve-outs (\textit{Mead} Step Zero, major questions, non-administered statutes) that left \textit{Chevron} a ``decaying husk'' unused by the Supreme Court since 2016; and (3) \textbf{Reliance}---rather than fostering stability, \textit{Chevron} destroyed reliance by inviting agencies to flip-flop across administrations without congressional action (\textit{id.} at 407--412).

\section{Concurring and Dissenting Opinions}

\begin{itemize}
  \item \textbf{Justice Thomas (Concurring):} Joined the opinion in full but wrote separately to argue that \textit{Chevron} also violated the Constitution's separation of powers. Regardless of what the APA says, mandatory deference strips courts of Article III ``judicial Power'' while transferring either judicial power or Article I ``legislative Powers'' to Article II executive agencies (\href{https://www.law.cornell.edu/supremecourt/text/22-451}{603 U.S. at 413--416}).
  \item \textbf{Justice Gorsuch (Concurring):} Declared that ``the Court places a tombstone on \textit{Chevron} no one can miss'' (\textit{id.} at 416). Tracing common-law \textit{stare decisis} from Blackstone and Hale to Lincoln, Gorsuch argued that \textit{Chevron} defied the APA, tilted the scales of justice systematically toward the government, undermined due process fair notice and the rule of lenity, and inflicted ``regulatory whiplash'' on ordinary citizens like veterans and immigrants (\textit{id.} at 416--448).
  \item \textbf{Justice Kagan (Dissenting, joined by Sotomayor, and by Jackson in \textit{Relentless}):} Defended \textit{Chevron} as a sensible, forty-year default rule reflecting Congress's likely preference that expert, experienced, politically accountable agencies---not generalist judges---resolve technical and policy-laden statutory gaps (e.g., whether an alpha amino acid polymer is a ``protein'' under FDA law or which squirrel populations are ``distinct'' under the Endangered Species Act). Kagan argued APA \S~706 is indeterminate (mentioning neither deference nor \textit{de novo} review for legal questions), noted Congress repeatedly legislated against \textit{Chevron}'s backdrop without repealing it, and condemned the majority for replacing ``a rule of judicial humility'' with ``a rule of judicial hubris'' (\textit{id.} at 448--479).
\end{itemize}

\section{Key Quotations}

\begin{quote}
``\textit{Chevron} is overruled. Courts must exercise their independent judgment in deciding whether an agency has acted within its statutory authority, as the APA requires. Careful attention to the judgment of the Executive Branch may help inform that inquiry. And when a particular statute delegates authority to an agency consistent with constitutional limits, courts must respect the delegation, while ensuring that the agency acts within it. But courts need not and under the APA may not defer to an agency interpretation of the law simply because a statute is ambiguous.''\\
--- Chief Justice Roberts, \href{https://www.law.cornell.edu/supremecourt/text/22-451}{603 U.S. at 412--413}
\end{quote}

\begin{quote}
``In an agency case as in any other, though, even if some judges might (or might not) consider the statute ambiguous, there is a best reading all the same---`the reading the court would have reached' if no agency were involved. \textit{Chevron}, 467 U.S., at 843, n. 11. It therefore makes no sense to speak of a `permissible' interpretation that is not the one the court, after applying all relevant interpretive tools, concludes is best. In the business of statutory interpretation, if it is not the best, it is not permissible.''\\
--- Chief Justice Roberts, \href{https://www.law.cornell.edu/supremecourt/text/22-451}{603 U.S. at 400}
\end{quote}

\begin{quote}
``By doing so, however, we do not call into question prior cases that relied on the \textit{Chevron} framework. The holdings of those cases that specific agency actions are lawful---including the Clean Air Act holding of \textit{Chevron} itself---are still subject to statutory stare decisis despite our change in interpretive methodology.''\\
--- Chief Justice Roberts, \href{https://www.law.cornell.edu/supremecourt/text/22-451}{603 U.S. at 412}
\end{quote}

\begin{quote}
``A rule of judicial humility gives way to a rule of judicial hubris. ... In one fell swoop, the majority today gives itself exclusive power over every open issue---no matter how expertise-driven or policy-laden---involving the meaning of regulatory law. As if it did not have enough on its plate, the majority turns itself into the country's administrative czar.''\\
--- Justice Kagan, dissenting, \href{https://www.law.cornell.edu/supremecourt/text/22-451}{603 U.S. at 449--450}
\end{quote}

\section{Doctrinal Significance and Modern Legacy (Through 2026)}

\textit{Loper Bright} reshaped American administrative law and separation-of-powers litigation across federal courts through 2026:

\begin{itemize}
  \item \textbf{Amplification via \textit{Corner Post} and Circuit Invalidations:} Decided three days later, \href{https://www.supremecourt.gov/opinions/23pdf/22-1008_1b82.pdf}{\textit{Corner Post, Inc. v. Board of Governors}}, 603 U.S. 799 (2024), held that the six-year APA limitations period begins when a plaintiff is injured rather than when a rule is promulgated, exposing older regulations to fresh \textit{Loper Bright} review. Applying \textit{Loper Bright} without \textit{Chevron} deference, the Sixth Circuit in \href{https://www.law.cornell.edu/wex/net_neutrality}{\textit{Ohio Telecom Ass'n v. FCC}} (\textit{In re MCP No. 185}), 124 F.4th 993 (6th Cir. 2025), vacated the FCC's 2024 Title II net-neutrality reclassification of broadband internet service.
  \item \textbf{Clarified Boundaries in 2025--2026 Supreme Court Decisions:} The Court has simultaneously checked overbroad readings of \textit{Loper Bright}. In \href{https://www.supremecourt.gov/opinions/24pdf/23-975_m648.pdf}{\textit{Seven County Infrastructure Coalition v. Eagle County}}, 605 U.S. \_\_\_ (2025), the Court distinguished pure statutory interpretation under \textit{Loper Bright} from substantial judicial deference owed to agency technical and factual judgment in NEPA environmental reviews. Likewise, in \href{https://supreme.justia.com/cases/federal/us/607/24-777/}{\textit{Urias-Orellana v. Bondi}}, 607 U.S. \_\_\_ (2026), the Court confirmed that \textit{Loper Bright} does not displace express statutory deferential review standards (such as INA \S~1252(b)(4)(B)) or lawful delegations of agency authority.
  \item \textbf{Rise of ``\textit{Skidmore}'' and Express-Delegation Battles:} Through 2026, agency litigation has pivoted sharply toward two fronts: persuading courts under \textit{Skidmore}'s thoroughness-and-consistency factors and proving that statutory terms like ``reasonable,'' ``appropriate,'' or ``as defined by regulation'' constitute express \textit{Loper Bright} delegations of policymaking discretion.
\end{itemize}

\section{Pedagogical Critiques and Exam Hypotheticals}

\subsection{Critiques}

\begin{itemize}
  \item \textbf{The ``Single Best Meaning'' Fiction:} Critics echo Justice Kagan's argument that many technical regulatory phrases genuinely lack a single textually fixed meaning at enactment, meaning \textit{Loper Bright} transfers unavoidable policy choices from accountable agencies to life-tenured judges---risking circuit splits and ideological divergence.
  \item \textbf{Boundary Uncertainty Over ``Delegations'':} Because Roberts preserved deference where a statute delegates discretion (e.g., ``appropriate'' or ``reasonable''), scholars note that disputes once fought at \textit{Chevron} Step One have migrated into an equally contested battle over which statutory terms qualify as express delegations.
  \item \textbf{Tension with Statutory \textit{Stare Decisis}:} Lower courts and litigants struggle to reconcile \textit{Loper Bright}'s assurance that prior \textit{Chevron} holdings remain intact with agencies' attempts to revise old rules---since an agency can no longer rely on \textit{Brand X} to switch from one ``permissible'' reading to another.
\end{itemize}

\subsection{``Opposite Side'' Contrast: \textit{Chevron} and \textit{Seven County Infrastructure Coalition}}

Contrast \textit{Loper Bright} with \href{https://supreme.justia.com/cases/federal/us/467/837/}{\textit{Chevron U.S.A. Inc. v. NRDC}}, 467 U.S. 837 (1984), and \href{https://www.supremecourt.gov/opinions/24pdf/23-975_m648.pdf}{\textit{Seven County Infrastructure Coalition v. Eagle County}}, 605 U.S. \_\_\_ (2025). In \textit{Chevron}, statutory silence on whether ``stationary source'' meant a single machine or an entire plant triggered mandatory judicial deference to the EPA's policy accommodation. Under \textit{Loper Bright}, that same silence triggers independent judicial construction to find the single best reading. Yet \textit{Seven County} (and \textit{Urias-Orellana}, 2026) illustrates the opposite side of \textit{Loper Bright} itself: once a court independently confirms the statute's boundaries, agencies still receive substantial deference on factual findings, technical predictions, and reasoned policy choices within expressly delegated discretion.

\subsection{Exam Hypotheticals}

\begin{enumerate}
  \item \textbf{Express Delegation vs. Statutory Silence:} Congress empowers the EPA to set emissions limits for ``emerging synthetic compounds'' and directs the agency to promulgate rules ``as the Administrator deems reasonable and necessary to protect public health.'' EPA (a) defines ``emerging synthetic compound'' to include a legacy chemical used since 1950, and (b) chooses a 10-parts-per-billion threshold after weighing scientific studies and compliance costs. Under \textit{Loper Bright}, a court reviews (a) \textit{de novo} using traditional statutory tools (granting at most \textit{Skidmore} persuasive weight), whereas under (b) the court recognizes Congress's express ``reasonable and necessary'' delegation and reviews EPA's threshold deferentially for reasoned decisionmaking under APA \S~706(2)(A).
  \item \textbf{Post-\textit{Loper Bright} Agency Flip-Flop:} In 2012, a court of appeals upheld an FDA rule under \textit{Chevron} Step Two. In 2026, a new administration repeals that rule and adopts the opposite statutory interpretation. Under \textit{Loper Bright} (603 U.S. at 412), statutory \textit{stare decisis} protected the 2012 holding that the original rule was lawful---but without \textit{Chevron} and \textit{Brand X}, can the agency defend its new interpretation as merely ``permissible,'' or must the reviewing court now independently determine the single best meaning of the statute?
\end{enumerate}

\end{document}
